% !TeX root = ../../relatorio-projeto-aplicado-exemplo-undf.tex
\undfsetcontent{resumo}{%
Este relatório apresenta um exemplo didático de Projeto Aplicado sobre
validação de solicitações de reserva de laboratório. O objetivo foi especificar
regras simples, implementar um validador e conferir os resultados em entradas
sintéticas. O contexto acadêmico, o semestre de referência e os integrantes
da equipe são fictícios; o programa e os casos acompanham o exemplo e podem
ser executados localmente. O método compreendeu definição de requisitos,
delimitação de escopo, implementação em Python e comparação entre retornos
obtidos e resultados previamente esperados. Foram examinados campos
obrigatórios, interpretação de horários, ordem do intervalo e limite de
capacidade. Os seis casos produziram os retornos esperados, com duas entradas
aceitas e quatro recusadas por impedimentos distintos. O resultado demonstra
o comportamento observado nesse conjunto, sem comprovar ausência de defeitos
ou adequação a uma rotina institucional. Interface, armazenamento e conflitos
entre reservas ficaram fora da entrega. Não houve pesquisa com pessoas,
implantação em laboratório ou avaliação de impacto operacional. Concluiu-se
que o recorte técnico foi concluído com evidências reproduzíveis e limitações
explícitas. O documento demonstra uma entrega final semestral e sua possível
rastreabilidade em um futuro portfólio, cuja consolidação permanece uma etapa
separada e dependente das orientações do curso.}
\undfsetcontent{palavraschave}{projeto aplicado; validação; relatório final}
\undfsetcontent{abstract}{%
This report presents an educational applied project concerning the validation
of laboratory reservation requests. The objective was to specify simple
rules, implement a validator and check its outputs using synthetic inputs.
The academic setting, reference semester and team members are fictional;
the program and cases are provided with the example and can be executed
locally. The method included requirements definition, scope delimitation,
implementation in Python and comparison between actual outputs and previously
defined expectations. Required fields, time interpretation, interval order
and capacity limits were examined. All six cases produced the expected
outputs, with two accepted inputs and four rejected inputs for distinct
reasons. These findings describe the behavior observed within this limited
set; they do not demonstrate the absence of defects or suitability for
institutional use. Interface development, storage and conflicts between
reservations were excluded from the delivery. No research with people,
laboratory deployment or operational impact assessment was conducted. The
technical scope was completed with reproducible evidence and explicit
limitations. The document illustrates a final semester report and the
traceability that could support its selection for a future portfolio.
Portfolio consolidation remains a separate step subject to degree program
requirements.}
\undfsetcontent{keywords}{applied project; validation; final report}
